\setlength{\baselineskip}{20pt}
\begin{abstract}
铁路道岔尖轨、基本轨等关键部件在长期服役中易产生表面及近表面伤损，威胁行车安全。现有检测方法在空间受限的道岔区域面临安装困难、抗干扰能力不足等局限。电磁层析成像（Electromagnetic Tomography, EMT）是一种基于交变电磁场感应效应的非接触式无损检测技术，具有可平面阵列化、适于近表层缺陷探测等特点，在道岔伤损检测中具有应用潜力。然而，道岔钢轨属于高电导率、高磁导率铁磁性目标，缺陷引起的电磁扰动微弱，且EMT逆问题具有显著病态性和非线性，传统基于灵敏度矩阵的线性重建方法难以兼顾检测精度、鲁棒性与工程适用性。

本文围绕便携式EMT硬件平台与智能成像算法两条主线展开研究，并构建从电磁响应采集到伤损成像显示的完整验证链路。在算法方面，针对经典深度算子网络（Deep Operator Network, DeepONet）Branch--Trunk结构在EMT多激励结构化输入与连续域输出适配上的不足，提出了一种共享条件特征驱动深度算子网络（Shared-Conditional-Feature Deep Operator Network, SCF-DeepONet）。该网络通过共享观测编码模块将多视角复数差分电压响应映射为全局条件特征、激励条件特征和二维局部特征图，并在此基础上并行构建二维伤损热图预测主分支与三维磁矢势差场重构辅助分支。网络架构消融实验表明，在相同训练条件下，SCF-DeepONet的平均交并比（Intersection over Union, IoU）相较原始DeepONet基线由67.79\%提升至81.53\%，困难样本检测下限（Worst-10 Avg IoU）由37.09\%提升至51.34\%，验证了所提出架构在铁磁性EMT伤损检测中的有效性。在此基础上，本文进一步将三维差场辅助监督作为训练阶段的物理感知正则化与物理信息神经网络（PINN）渐进预实验：设计了由热图预训练、热图--场量联合训练、热图精修和偏微分方程（Partial Differential Equation, PDE）弱物理微调构成的四阶段课表训练策略；实验表明，辅助监督虽未在同分布样本上取得更高平均精度，但使模型在圆形伤损零样本测试中的平均IoU由34.05\%提升至40.31\%，改善了分布外泛化能力，也为后续严格PDE约束的嵌入提供了经验基础与初始化前提。

在硬件方面，设计并实现了一套便携式EMT检测平台。系统采用“测控板+信号切换板”的双板式架构，构建了$4 \times 4$平面线圈阵列传感器、收发异构大电流多路选通网络、基于直接数字频率合成（Direct Digital Synthesis, DDS）的激励产生与双级功率放大链路、双相锁相解调与高精度模数转换模块，并完成了主控通信、电源管理及PCB抗干扰设计，可在100 kHz主工作点下稳定获取多通道复数互感响应。系统单帧完整测量周期约为1.46 s，基础性能测试表明激励信号稳定、锁相解调链路工作正常、通道一致性良好。

最后，通过机加工仿损试块实验，验证了所研制的便携式硬件平台与SCF-DeepONet算法联合工作的可行性。实验结果表明，SCF-DeepONet对缺陷位置的响应较传统线性成像方法更加集中，所构建的“硬件采集--算法反演”端到端链路具备一定的工程应用潜力。
		
\noindent\keywords{电磁层析成像；无损检测；深度算子网络；多任务学习；物理信息辅助学习}
\end{abstract}